\icode{O-A-SU-260929-001}

	\lkf{1.天然放射现象}

	1896年法国贝克勒尔发现天然放射现象

	\begin{table}[htbp]
		\centering
		\begin{tblr}{c|c|c|c}
		\hline
		&\(\alpha\) 射线&\(\beta\)射线&\(\gamma\)射线\\
		\hline
		组成&高速氦核（\ce{^4_2He}）流&高速电子（\ce{^0_{-1}e}）流&光子流（高频电磁波）\\
		\hline
		电荷&\( + 2e\)&\( - e\)&0\\
		\hline
		速度&\(0.1c\) &\(0.99c\)&\(c\)\\
		\hline
		贯穿本领&最弱&较弱&最强\\
		\hline
		电离能力&很强&较弱&最弱\\
		\hline
		\end{tblr}
		\end{table}

		元素的放射性不受其存在形式是单质还量化合物的影响，也不受外界环境因素（如温度，压强等）变化的影响

	\lkf{2.天然衰变}

	\lkf{\(\alpha\)衰变：}原子核放出\(\alpha\)粒子(\ce{^4_2He})的衰变，特征：一变二+有氦核+\(\gamma\)射线

	\qquad\ce{^A_ZX \to ^{A-4}_{Z-2}Y + ^4_2He}

	\lkf{\(\beta\)衰变：}原子核放出\(\beta\)粒子(\ce{^0_{-1}e})的衰变，特征：一变二+有电子+\(\gamma\)射线

	\qquad\ce{^A_ZX \to ^{A}_{Z+1}Y + ^0_{-1}e}

	\lkf{半衰期：}有一半原子核发生衰变所用的时间，和核外因素无关\(m = m_0(\frac{1}{2})^{\frac{t}{T}} = m_0(\frac{1}{2})^n\)

	\lkf{3.重核裂变}

	\lkf{核裂变：}用中子(\ce{^1_0n})轰击重核，使之分裂为多个核并放出多个中子

	\qquad\ce{^{235}_{92}U + ^1_0n \to ^{144}_{56}Ba + ^{89}_{36}Kr + 3 ^1_0n}

	\lkf{4.轻核聚变}

	\lkf{核聚变：}多个轻核在超高温下转变为重核，并放出中子

	\qquad\ce{^2_1H + ^3_1H \to ^4_2He + ^1_0n}

	\lkf{5.人工核反应}

	\lkf{定义：}用\ce{^4_2He}/\ce{^1_1H}/\ce{^1_0n}等轰击旧核使之变为新核，也称核反应

	卢瑟福发现质子：\ce{^4_2He + ^{14}_{7}N \to ^{17}_{8}O + ^1_1H}

	查德威克发现中子：\ce{^4_2He + ^9_4Be \to ^{12}_{6}C + ^1_0n}

	约里奥\(\cdot\)居里夫妇发现放射性同位素，同时发现正电子：

	\ce{^4_2He + ^{27}_{13}Al \to ^{30}_{15}P + ^1_0n}\qquad\ce{^{30}_{15}P \to ^{30}_{14}Si + ^0_{+1}e}

	\lkf{6.核反应中的守恒}
	
	所有核反应都满足“质量数守恒”和“电荷数守恒”。

	\lkf{7.质能方程}

	\(\Delta E = \Delta mc^2\)

	\lkf{8.结合能与比结合能}

	原子核比结合能越大越稳定

	所有的核反应都是向比结能大的方向进行，即生成物比反应物稳定